
Este circuito se implementa conectando la neurona electrónica del laboratorio del \textit{GNB} \cite{Szcs2000} a la \textit{DAQ Board} (Figura \ref{FIG:MODELELECTROSYN}) y luego conectando de forma bidireccional, con sinapsis eléctrica, el modelo de Hindmarsh-Rose \cite{Hindmarsh1984, Golowasch1999} mediante la \textit{DAQ Board} y viceversa. Si la corriente de la sinapsis es negativa (este es un parámetro ajustable en el plugin), la actividad de las neuronas estará en antifase, esto significa que será en forma de secuencia bifásica alternante, es decir, cada neurona se activa después de que termina la inhibición de la otra (Figura \ref{FIG:HRELECTBISYNANTI}). Por otro lado, si la corriente es positiva, la actividad estará en fase y ambas neuronas se sincronizarán (Figura \ref{FIG:HRELECTBISYNFASE}), ya que por la sinapsis eléctrica fluye corriente para igualar los potenciales de las neuronas. Este comportamiento entre modelo y neurona electrónica, sugiere que se podrían obtener unos resultados análogos con una neurona viva \cite{Pinto2000}, dado que las figuras mencionadas validan la eficacia de la implementación en tiempo real para dichas interacciones. La comunicación entre los plugins que implementan el ciclo-cerrado de la conexión se detalla en la Figura \ref{FIG:PLUGHREHR}.

\begin{figure}[Esquema de la conexión entre el modelo y la neurona electrónica]{FIG:MODELELECTROSYN}{Conexíon entre el modelo de Hindmarsh-Rose en ejecución en el ordenador y el de la neurona electrónica.}
    \image{0.3\textwidth}{}{model-electronic}
\end{figure}\textbf{}

\begin{figure}[Conexión bidireccional modelo-electrónica en antifase]{FIG:HRELECTBISYNANTI}{Conexión en antifase del modelo Hindmarsh-Rose (\textit{model\_neuron}) y de la neurona electrónica (\textit{living\_neuron}).}
    \image{0.9\textwidth}{}{v_hr-e_hr-syn_antiphase}
\end{figure}

\begin{figure}[Conexión bidireccional modelo-electrónica en fase]{FIG:HRELECTBISYNFASE}{Conexión en fase del modelo Hindmarsh-Rose (\textit{model\_neuron}) y de la neurona electrónica (\textit{living\_neuron}).}
    \image{0.9\textwidth}{}{v_hr-e_hr-syn_phase}
\end{figure}

\begin{figure}[Esquema de conexión de plugins para crear el circuito entre el modelo y la neurona electrónica]{FIG:PLUGHREHR}{Conexiones de plugins realizada para establecer la connectividad entre el modelo de Hindmarsh-Rose en ejecución en el ordenador y el de la neurona electrónica.}
    \image{0.5\textwidth}{}{hr-ehr}
\end{figure}\textbf{}

Este circuito se ha probado también enviando la señal del modelo neuronal a través de la \textit{DAQ Board}, debido a la problemática de \textit{RTXI} con la conexión de los plugins en ciclo que se ha comentado con anterioridad. Para ello, se ha realizado un puente entre una de las salidas de la tarjeta y se ha conectado a una de las entradas. Los parámetros usados en esta ejecución han sido distintos, pero igualmente se consigue tanto la interacción en fase (Figura \ref{FIG:HRELECTBISYNANTIFORD}) como en antifase (Figura \ref{FIG:HRELECTBISYNFASEFORD}). Las conexiones de los plugins se detallan en la Figura \ref{FIG:PLUGHREHRDAQ}.

\begin{figure}[Sinapsis bidireccional modelo-electrónica en antifase con forwarding]{FIG:HRELECTBISYNANTIFORD}{Sinapsis en antifase del modelo Hindmarsh-Rose (\textit{model\_neuron}) y de la neurona electrónica (\textit{living\_neuron}) pasando la conexión por la \textit{DAQ Board}.}
    \image{0.9\textwidth}{}{v_hr-e_hr-daq_forwarding-syn_antiphase}
\end{figure}

\begin{figure}[Sinapsis bidireccional modelo-electrónica en fase con forwarding]{FIG:HRELECTBISYNFASEFORD}{Sinapsis en fase del modelo Hindmarsh-Rose (\textit{model\_neuron}) y de la neurona electrónica (\textit{living\_neuron}) pasando la conexión por la \textit{DAQ Board}.}
    \image{0.9\textwidth}{}{v_hr-e_hr-daq_forwarding-syn_phase}
\end{figure}

\begin{figure}[Esquema de conexión de plugins para la sinapsis entre el modelo y la neurona electrónica con DAQ forwarding]{FIG:PLUGHREHRDAQ}{Conexiones de plugins realizada para crear el circuito entre el modelo de Hindmarsh-Rose en ejecución en el ordenador y el de la neurona electrónica pasando la salida del modelo por la \textit{DAQ Board} para evitar el ciclo.}
    \image{0.5\textwidth}{}{hr-ehr-daq}
\end{figure}\textbf{}
